\section{Extra Content}
This content is not required for the AP Chemistry exam, but is relevant to the study of chemistry and may contain helpful information.
\subsection{Nuclear Chemistry}
Nuclear chemistry is the study of the nucleus of the atom and the reactions that occur within it. This includes topics such as radioactive decay, nuclear fission, and nuclear fusion. 
\subsubsection{Radiation}
Radiation is the emission of energy in the form of particles or electromagnetic waves. All objects have some for of electromagnetic radiation, but only some have nuclear radiation. Nuclear radiation is emitted by unstable nuclei as they decay into more stable forms.
The radiation that we receive from the sun and space is called the cosmic background radiation. This radiation exists everywhere and is a normal part of our environment. In addition, we also use radiation such as X-rays in medicine. 

Radioisotopes are unstable isotopes that emit radiation as they decay. They can release two types of radiation, ionizing and non-ionizing. Ionizing radiation has enough energy to remove electrons from atoms, while non-ionizing radiation does not. Ionizing radiation can be harmful to living organisms, while non-ionizing radiation is generally considered safe.

The dangers of radiation depend on the type and amount of radiation, as well as the duration of exposure. High levels of ionizing radiation can cause damage to cells and DNA, leading to an increased risk of cancer and other health problems. To truly know the consequences of radiation, we need to consider how much, for how long, and how often we are exposed to it.
\subsubsection{Radioactive Decay}
Radioactive decay is the process by which an unstable atomic nucleus loses energy by emitting radiation. In Rutherford's lead box experiment, he found three types of radiation: alpha, beta, and gamma. Alpha particles ($^4_2\alpha$) are helium nuclei (\ce{^4_2He^2+}), and are the largest and least penetrating type of radiation. They have a +2 charge and can be stopped by a sheet of paper. A nucleus will undergo alpha radiation if there are too many protons.Beta particles ($^0_{-1}\beta$) are high-energy electrons (\ce{e^-}) that are emitted from the nucleus when a neutron decays into a proton. They have a -1 charge and can be stopped by a sheet of aluminum. A nucleus will undergo beta radiation if there are too many neutrons. Gamma rays ($^0_0\gamma$) are high-energy electromagnetic waves that are emitted from the nucleus when it undergoes a transition to a lower energy state. They have no charge and can be stopped by a thick layer of lead. A nucleus will undergo gamma radiation if it has excess energy after undergoing alpha or beta decay. A nucleus can also undergo a neutron emission ($^1_0n$), which is the emission of a neutron from the nucleus. 
\begin{table*}[ht]
    \centering
    \begin{tabular}{|c|c|c|c|c|c|c|}
        \hline
        \textbf{Type of Radiation} & \textbf{Symbol} & \textbf{Size} & \textbf{Charge} & \textbf{Speed} & \textbf{Penetrating Power} & \textbf{Identity} \\
        \hline
        Alpha & $^4_2\alpha$ & Large & +2 & Slow & Low & \ce{^4_2He^2+} \\
        Beta & $^0_{-1}\beta$ & Small & -1 & Fast & Medium & \ce{e^-} \\
        Gamma & $^0_0\gamma$ & None & 0 & Very Fast & High & \ce{\gamma} \\
        Neutron & $^1_0n$ & Medium & 0 & Fast & Medium & \ce{n} \\
        \hline
    \end{tabular}
    \caption{Types of Nuclear Radiation}
\end{table*}
\subsubsection{Nuclear Reactions}
Nuclear reactions involve changes in the nucleus of an atom, resulting in the formation of new elements or isotopes. When writing out a nuclear reaction, we need to ensure that the total mass number and atomic number are conserved. For example, in the beta decay of carbon-14, we have:
\[\ce{^{14}_6C -> ^{14}_7N + ^0_{-1}\beta}\]
In this reaction, the mass number is conserved (14 on both sides) and the atomic number is conserved (6 on the left and 7 - 1 = 6 on the right).

There are two main types of nuclear reactions: fission and fusion. Nuclear fission is the process by which a heavy nucleus splits into two or more smaller nuclei, releasing a large amount of energy. This process is used in nuclear power plants and atomic bombs. Nuclear fusion is the process by which two light nuclei combine to form a heavier nucleus, releasing an even larger amount of energy. This process powers the sun and other stars. The reactions in bombs are so powerful because they involve a chain reaction, where the products of one reaction can trigger additional reactions, leading to an exponential increase in energy release. 
\subsubsection{Half-Life}
The half-life of a radioactive isotope is the time it takes for half of the atoms in a sample to decay. The speed of this process is a first order reaction, meaning that the rate of decay is proportional to the number of atoms remaining. The half-life can be calculated using the same formula for any other first order reaction.
\subsection{Organic Chemistry}
Organic chemistry is the study of carbon-containing compounds and their properties. This includes topics such as hydrocarbons, functional groups, and reactions. This is the largest branch of chemistry in the modern world. It is involved in the production of pharmaceuticals, plastics, synthetic materials, solvents, fuels, and much more.

The reason that carbon is so important in organic chemistry is because of its unique ability to form four covalent bonds with other atoms, including itself. This allows for the formation of a wide variety of complex molecules with different shapes and properties. 
\subsubsection{Condensed Structural Formulas}
Condensed structural formulas are a way to represent the structure of a molecule in a more compact form than a full structural formula. In a condensed structural formula, the atoms are listed in a specific order, and the bonds between them are implied rather than explicitly shown. For example, the condensed structural formula for ethanol is 
\[\chemfig{CH_3 - CH_2 - OH},\]
which indicates that there are two carbon atoms connected by a single bond, with three hydrogen atoms attached to the first carbon and two hydrogen atoms and a bond to an oxygen atom attached to the second carbon. Condensed structural formulas are often used in organic chemistry to quickly convey the structure of a molecule without having to draw out the full structural formula.
\subsubsection{Alkanes}
Alkanes are the simplest type of hydrocarbon, consisting of only carbon and hydrogen atoms connected by single bonds. They have the general formula \ce{C_nH_{2n+2}}. The prefixes for naming alkanes are based on the number of carbon atoms in the longest continuous chain. The prefixes are as follows: 1 carbon - meth-, 2 carbons - eth-, 3 carbons - prop-, 4 carbons - but-, 5 carbons - pent-, 6 carbons - hex-, 7 carbons - hept-, 8 carbons - oct-, 9 carbons - non-, and 10 carbons - dec-, 11 carbons - undec-, 12 carbons - dodec-. For example, the alkane with 5 carbon atoms is called pentane (\ce{C5H12}). 

The longest continuous chain of carbon atoms in a molecule is called the stem or chain and is used to determine the base name of the compound. The other carbon atoms that are not part of the longest chain are called the branches or side chains. These side chains have the same prefixes as the alkanes, but with the suffix -yl instead of -ane. For example, a methyl group 
\[\chemfig{CH_3 - R}\] is a side chain with one carbon atom, and an ethyl group
\[\chemfig{CH_3 - CH_2 - R}\] is a side chain with two carbon atoms, where R represents the rest of the molecule.

When writing the name of an alkane, we first identify the longest continuous chain of carbon atoms and use the appropriate prefix to name it. Then, we identify any side chains and use their prefixes with the suffix -yl to name them. Finally, we combine the names of the main chain and the side chains, using numbers to indicate the position of each side chain on the main chain and ordering the side chains alphabetically. For example, the compound with the condensed structural formula 
\[\chemfig{CH_3-CH(-[:90]CH_3)-CH_2-CH_3}\] 
is named 2-methylbutane, because the longest chain has 4 carbon atoms (butane) and there is a methyl group attached to the second carbon atom of the main chain. If there are multiple side chains, we use prefixes such as di-, tri-, tetra-, etc. to indicate the number of identical side chains. The structural formula for 2,3-dimethylpentane is
\[\chemfig{CH_3-CH(-[:90]CH_3)-CH(-[:90]CH_3)-CH_2-CH_3}\]

The numbers used to indicate the position of side chains are always the lowest possible numbers. For example, if there is a choice between numbering the main chain from left to right or from right to left, we choose the direction that gives the side chains the lowest possible numbers. If there are multiple side chains, we number the main chain in the direction that gives the first side chain the lowest possible number. If there is a tie, we look at the second side chain, and so on until we can break the tie. 

When the chian creates a ring, we use the prefix cyclo- before the name of the alkane. For example, the compound with the condensed structural formula
\[\chemfig{CH_2*6(-CH_2-CH_2-CH_2-CH_2-CH_2-[,,2])}\]
is named cyclohexane. 
\subsubsection{Isomers}
Isomers are compounds that have the same molecular formula but different structural formulas. Take for example the molecular formula \ce{C4H10}. This formula corresponds to two different compounds: butane and 2-methylpropane. Butane has a straight chain of four carbon atoms, while 2-methylpropane has a branched chain with three carbon atoms in the main chain and one carbon atom as a side chain. These two compounds have different physical and chemical properties, even though they have the same
 molecular formula.
\subsubsection{Functional Groups}
Functional groups are specific patterns of elements and bonds that give a molecule its characteristic properties and reactivity. Compounds can combine several of these groups to make more complex molecules.

To name a compound with a functional group, we first identify the longest continuous chain of carbon atoms that contains the functional group and use the appropriate prefix to name it. Then, we identify the functional group and use its name as a suffix to indicate its presence in the molecule. 

\begin{table*}
    \centering
    {\setlength{\extrarowheight}{20pt}
    \begin{tabular}{|c|c|c|c|}
        \hline
        \textbf{Name} & \textbf{Functional Group} & \textbf{Example} & \textbf{Prefix/Suffix} \\
        \hline 
        Alkanes & $\chemfig{R-C-C-R}$ & \ce{C_4H_{10}} (Butane) & -ane \\
        Alkenes & $\chemfig{R-C=C-R}$ & \ce{C_4H_8} (2-Butene) & -ene \\
        Alkynes & $\chemfig{R-C~C-R}$ & \ce{C_4H_6} (2-Butyne) & -yne \\
        Alcohols & $\chemfig{R-OH}$ & \ce{C_2H_5OH} (Ethanol) & -anol \\
        Acids & $\chemfig{R-C(=[:90]O)-OH}$ & \ce{CH_3COOH} (Ethanoic Acid) & -oic acid \\
        Aldehydes & $\chemfig{R-C(=[:90]O)-H}$ & \ce{CH_3CHO} (Ethanal) & -anal \\
        Ketones & $\chemfig{R-C(=[:90]O)-R}$ & \ce{CH_3COCH_3} (Propanone) & -one \\
        Esters & $\chemfig{R-C(=[:90]O)-O-R}$ & \ce{C_3COOC_3} (Methyl Ethanoate) & -anoate \\
        Ethers & $\chemfig{R-O-R}$ & \ce{C_2H_5OC_2H_5} (Diethyl Ether) & Ether \\
        Rings & $\chemfig{CH_2*6(-CH_2-CH_2-CH_2-CH_2-CH_2-[,,2])}$ & \ce{C_6H_{12}} (Cyclohexane) & cyclo- \\
        Benzene & $\chemfig{{CH}*6(-CH=CH-CH=CH-CH=[,,2,2])}$ & \ce{C_6H_6} (Benzene) & Benzene \\
        \hline
    \end{tabular}}
    \caption{Common Functional Groups in Organic Chemistry}
    \label{tab:functional_groups}
\end{table*}

The most common functional groups in organic chemistry are alkanes, alkenes, alkynes, alcohols, acids, aldehydes, ketones, esters, ethers, rings, and benzene. Each of these functional groups has its own unique properties and reactivity, which can be used to predict the behavior of the molecule in chemical reactions. For now just worry about the naming of these functional groups, but in the future you may want to learn about their properties and reactions as well. There are many more than the functional groups listed in Table \ref{tab:functional_groups}, but these are some of the most common ones that you will encounter in organic chemistry.